%% kapa: exam=1 round=33 year=2022 period=2 form=A subject=appraisal_law subject_ko=감정평가관계법규 seq=26 number=26
%% source: Q-Net 감정평가사 기출문제 (cst00309) artlSeq=5209997 file=raw/1008/5209997/2022년도 제33회 감정평가사 1차시험 2교시(감정평가 관계법규, 회계학).pdf
%% conversion: kordoc 4.24.51 -> tools/gen1.py (auto); review=proofread-sonnet; flags=none
\begin{question}{26}
건축법령상 건축물의 용도에 따른 건축허가의 승인에 관한 설명이다. ( )에 해당하는 건축물이 아닌 것은?
\begin{framedtext}
시장·군수가 자연환경이나 수질을 보호하기 위하여 도지사가 지정·공고한 구역에 건축하는 3층 이상 또는 연면적의 합계가 1천제곱미터 이상인 건축물로서 ( )의 건축을 허가하려면 미리 도지사의 승인을 받아야 한다. \par
\end{framedtext}
\begin{choices}
  \item 공동주택
  \item 제2종 근린생활시설(일반음식점만 해당한다)
  \item 업무시설(일반업무시설은 제외한다)
  \item 숙박시설
  \item 위락시설
\end{choices}
\end{question}
\answer{3}
